\section{Generated Data}
\label{sec:generator-eval}

\paragraph{Generator.} The generator samples a target path, assigns labels to
constants and predicates, then adds distractor branches that remain locally
coherent but do not derive the queried conclusion. In pseudocode:

\begin{verbatim}
sample target answer label
sample initial object state
for depth step in 1..D:
    sample active branch rule(s)
    derive next object state
    add coherent distractor rules for inactive branches
render one prompt
render the same latent proof as formal and English targets
\end{verbatim}

\paragraph{Uniqueness of the answer.}
Every metric compares a generated answer against one labelled answer, which
is only meaningful if the labelled answer is the sole one derivable from the
premises. Each example introduces fresh constants, so no constant is reused
across layers and no two branches can derive the same atom through reuse of
a state predicate. A preflight check computes the full Horn closure of
1{,}000 training instances and 2{,}000 evaluation instances with a solver
independent of the generator and records the number of candidate final
states derivable in each. The unique-answer rate is 1.0, the maximum number
of derivable candidates is 1, there are no generation failures, and answer
positions are balanced across the four candidates.

\paragraph{Trace grammar and checker.}
Formal traces contain four blocks, constants, predicates, premises and proof
lines, followed by a conclusion and an answer. A proof line contains a
formula and a rule justification such as premise retrieval or implication
elimination. The checker parses the generated blocks and verifies that each
proof line is licensed by earlier lines and premises. English traces use
controlled sentences for the same premises and proof steps, and a
deterministic translator maps those sentences back into the checker. The
formal surface is:

\begin{verbatim}
<formal>
<constants> CONST_DEF+ </constants>
<predicates> PRED_DEF+ </predicates>
<premises> FORMULA+ </premises>
<proof> PROOF_LINE+ </proof>
<conclusion> ATOM </conclusion>
<answer> LABEL </answer>

CONST_DEF  ::= CONST "=" LABEL
PRED_DEF   ::= PRED "x:" "x is" LABEL
ATOM       ::= PRED CONST
FORMULA    ::= ATOM | ATOM "&" ATOM "->" ATOM | ATOM "->" ATOM
PROOF_LINE ::= FORMULA "; R," REF
             | ATOM "; ->E," REF "," REF
             | ATOM "&" ATOM "; &I," REF "," REF
\end{verbatim}

References must point to earlier proof lines or listed premises, so a line
cannot justify a future step. Answer extraction searches for the final
\texttt{<answer>} field, and a generation without a well-formed final answer
is counted incorrect. Common checker failures are malformed block structure,
undeclared symbols, references to unavailable lines, applying implication
elimination to the wrong antecedent, and stopping before the queried
conclusion. English validation adds one further failure mode, in which the
controlled sentence must translate back into a formal step before the checker
can evaluate it.

\paragraph{A paired example.}
\label{sec:paired-example}
The prompt and the answer are identical within each pair. Training examples
contain more distractor premises and longer derivations.

\begin{verbatim}
Prompt:
a is cobalt. a is east.
If a is east and a is cobalt, then b is maple.
If b is maple, then b is south.
If b is south and b is maple, then c is cedar.
What is c?
\end{verbatim}

\noindent\begin{minipage}[t]{0.48\linewidth}
\textbf{\formal{} target}
\scriptsize
\begin{verbatim}
<formal>
<constants> a = a b = b c = c </constants>
<predicates> Cx: x is cobalt Ex: x is east
Mx: x is maple Sx: x is south Dx: x is cedar
</predicates>
<premises> Ca Ea
Ea & Ca -> Mb
Mb -> Sb
Sb & Mb -> Dc </premises>
<proof> Ca ; R,1
Ea ; R,2
Ea & Ca ; &I,7,6
Mb ; ->E,3,8
Sb ; ->E,4,9
Sb & Mb ; &I,10,9
Dc ; ->E,5,11 </proof>
<conclusion> Dc </conclusion>
<answer> cedar </answer>
\end{verbatim}
\end{minipage}\hfill
\begin{minipage}[t]{0.48\linewidth}
\textbf{\nat{} target}
\scriptsize
\begin{verbatim}
<think>
a is cobalt. a is east.
Since a is east and a is cobalt, b is maple.
Since b is maple, b is south.
Since b is south and b is maple, c is cedar.
Therefore, c is cedar.
</think>
<answer> cedar </answer>
\end{verbatim}
\end{minipage}

\section{Training Details}
\label{sec:sweep-details}

Every model trains for one epoch on $10^5$ examples, full-parameter under
fully sharded data parallelism with gradient checkpointing, in bfloat16,
with a global batch of 128 sequences, a maximum sequence length of 8{,}192,
warmup ratio 0.03 and learning rate $5\times10^{-6}$. The learning rate is
the one used by our earlier instruction-tuning runs on the same corpus and
is held fixed across conditions. The 3 runs per configuration on Qwen2.5-7B
differ only in the random seed, and only the final checkpoint is kept. Before
training starts, the mixture is tokenized and the examples whose target
would be cut by the sequence length are counted. The run aborts if that
count exceeds 0.5 percent, which prevents a condition from training on
derivations whose conclusions have been removed. No run aborted. Qwen2.5-32B
and the second-corpus runs use the same configuration, the former on two
nodes.

\section{Evaluation Details}
\label{sec:metrics}

\paragraph{Prompts and scoring.} ProofWriter and FOLIO are prompted with the
theory, the claim and the question whether the claim is true, false or
unknown (uncertain for FOLIO), and the answer is the first label in the
generation. The reasoning prompt for ProofWriter asks the model to reason
step by step and to give the final answer on its own last line. GPQA-Diamond
uses the four-option mirror of the benchmark and asks the model to work
through the options briefly and end with ``Answer: X''. BIG-Bench Hard ships
as a chain-of-thought few-shot task whose answer extractor reads a plain
continuation, and applying the chat template to its shots makes the
extractor match nothing for every condition, so this one suite is run without
the template. The LongBench multi-hop datasets use the standard prompt and
report F1. \branchproof{} prompts the model in the format of its training
documents and scores the fraction of items whose derivation the checker
accepts and whose answer is correct.

\paragraph{Benchmark subsets.} The chain subtasks of BIG-Bench Hard are web
of lies, tracking shuffled objects (three, five and seven objects), logical
deduction (three, five and seven objects) and formal fallacies. The
quantitative subset of GPQA-Diamond is defined by two regular expressions
applied to the question text. An item is quantitative if the question
contains a number followed by a physical unit (nm, mm, cm, km, kg, mol, ml,
eV, keV, MeV, GeV, K, Hz, MHz, GHz, degrees, percent, pc, AU, or a solar
mass or radius symbol), or if it contains a digit and asks for a value,
ratio, factor, energy, wavelength, number, probability, concentration, pH,
mass, radius, velocity or temperature, or contains ``calculate'',
``compute'', ``how many'', ``by what factor'' or ``what fraction''. The rule
assigns 73 items to the quantitative subset and 125 to the remainder, and it
depends only on the question text. The sign test compares, for each item,
the pooled runs trained with derivations against the pooled control runs.

\paragraph{Overlap audit.} No evaluation item of ProofWriter, FOLIO or the
BIG-Bench Hard subtasks shares an eight-gram with any training derivation.

\section{Additional Results}
\subsection{General benchmarks}
\label{sec:general-table}

Table~\ref{tab:general} reports nine general benchmarks. GSM8K is five-shot
with a chain of thought, HumanEval and MBPP report pass@1 from a single
greedy sample, and MMLU, ARC-Challenge, LogiQA, HellaSwag, PIQA and
WinoGrande are scored by log-likelihood over the answer options with the
default prompts of the evaluation harness. The 3 runs of the Qwen2.5-7B
control span 0.3 points on GSM8K, 0.1 on MMLU, 0.8 on ARC-Challenge, 0.6 on
LogiQA, 0.0 on HellaSwag, 0.4 on PIQA and 1.0 on WinoGrande. Rows for the
English rendering at 25 percent and some code rows are still being
evaluated.

\begin{table}[t]
\caption{General benchmarks for Qwen2.5-7B, in percent. GSM8K is five-shot
with a chain of thought, HumanEval and MBPP report pass@1, and the remaining
benchmarks are scored by log-likelihood over the answer options. Entries
marked -- are still being evaluated.}
\label{tab:general}
\begin{center}
\footnotesize
\begin{tabular}{llccccccccc}
\toprule
 & & GSM8K & MMLU & ARC-C & LogiQA & HellaSwag & PIQA & WinoGrande & HumanEval & MBPP \\
\midrule
\multicolumn{2}{l}{Control} & 77.0 & 70.3 & 49.6 & 36.8 & 73.2 & 80.5 & 71.0 & 66.8 & 65.2 \\
\midrule
\multirow{4}{*}{\formal} & 1\%  & 76.5 & 70.2 & 52.4 & 38.0 & 73.4 & 81.0 & 70.6 & 67.1 & 64.8 \\
 & 5\%  & 76.7 & 69.7 & 50.7 & 36.8 & 73.2 & 80.8 & 70.3 & -- & -- \\
 & 10\% & 76.2 & 70.0 & 50.4 & 38.9 & 73.2 & 80.9 & 71.0 & 68.3 & 65.2 \\
 & 25\% & 76.8 & 69.5 & 49.8 & 38.4 & 73.6 & 80.7 & 71.3 & -- & -- \\
\midrule
\multirow{4}{*}{\nat} & 1\%  & 76.5 & 70.2 & 51.1 & 37.5 & 73.3 & 80.6 & 71.3 & 65.9 & 65.2 \\
 & 5\%  & 75.7 & 70.2 & 51.9 & 36.6 & 73.6 & 81.0 & 70.2 & 63.4 & 64.6 \\
 & 10\% & 76.4 & 70.3 & 49.9 & 38.5 & 73.5 & 80.8 & 70.7 & 69.5 & 64.8 \\
 & 25\% & 78.2 & 69.9 & 50.7 & 37.6 & 73.5 & 80.5 & 71.4 & 65.2 & 66.4 \\
\bottomrule
\end{tabular}
\end{center}
\end{table}



\subsection{BIG-Bench Hard by subtask}
\label{sec:bbh-table}

\begin{table}[t]
\caption{BIG-Bench Hard resolved by subtask, in percent, mean over seeds.
``Chain'' pools the eight subtasks whose answer follows from propagating a
truth value or a state along a sequence of steps, which are web of lies, the
three object-tracking subtasks, the three logical-deduction subtasks and formal
fallacies. ``Others'' pools the remaining nineteen. The rightmost column gives
web of lies alone, the subtask closest in form to the generated data. The gain
is confined to the chain subtasks and grows with the replaced share.}
\label{tab:bbh}
\begin{center}
\begin{tabular}{llcccc}
\toprule
Notation & Share & macro & chain (8) & others (19) & web of lies \\
\midrule
none    & 0\%  & $65.0$ & $62.2$ & $66.1$ & $85.5$ \\
formal  & 1\%  & $66.2$ & $66.1$ & $66.3$ & $92.7$ \\
formal  & 5\%  & $67.2$ & $67.6$ & $67.1$ & $95.9$ \\
formal  & 10\% & $\mathbf{67.7}$ & $\mathbf{68.9}$ & $67.2$ & $\mathbf{97.6}$ \\
formal  & 25\% & $66.9$ & $67.6$ & $66.5$ & $98.0$ \\
English & 1\%  & $65.3$ & $63.2$ & $66.1$ & $84.1$ \\
English & 5\%  & $66.6$ & $65.2$ & $67.2$ & $84.5$ \\
English & 10\% & $66.9$ & $66.4$ & $67.2$ & $87.5$ \\
English & 25\% & $66.6$ & $65.6$ & $67.0$ & $87.6$ \\
\bottomrule
\end{tabular}
\end{center}
\end{table}


\subsection{Web of lies by gold answer}
\label{sec:wol-table}

Table~\ref{tab:wol} resolves web of lies by gold answer and by the number
of negating links, and Table~\ref{tab:steps} reads the intermediate value at
every position of the chain of thought.

\begin{table}[H]
\caption{Web of lies on Qwen2.5-7B, the BIG-Bench Hard subtask that chains
assertions about who tells the truth, accuracy in percent, mean over runs on
250 items. Gold answers are ``yes'' on 48.8 percent of items. A negating
link is a person who says that the previous person lies, and the column for
3 negating links covers the items with the most such links. Every condition
emits the answer phrase on every item and writes about 129 words.}
\label{tab:wol}
\begin{center}
\begin{tabular}{llccccc}
\toprule
Notation & Share & $P(\text{answers yes})$ & gold yes & gold no & 3 negating links & accuracy \\
\midrule
none    & 0\%  & $.567$ & $.932$ & $.781$ & $.717$ & $.855$ \\
formal  & 1\%  & $.500$ & $.937$ & $.917$ & $.888$ & $.927$ \\
formal  & 5\%  & $.484$ & $.954$ & $.964$ & $.957$ & $.959$ \\
formal  & 10\% & $.485$ & $.973$ & $\mathbf{.979}$ & $\mathbf{.973}$ & $\mathbf{.976}$ \\
English & 1\%  & $.569$ & $.921$ & $.766$ & $.690$ & $.841$ \\
English & 5\%  & $.571$ & $.926$ & $.768$ & $.690$ & $.845$ \\
English & 10\% & $.539$ & $.923$ & $.828$ & $.775$ & $.875$ \\
\bottomrule
\end{tabular}
\end{center}
\end{table}

\begin{table}[H]
\caption{Intermediate values in web of lies. The model writes 1 numbered line per person, so the truth value it writes can be read at each position and compared with the value the chain implies. The last two columns give the final answer accuracy when every intermediate value is correct and when at least one is not.}
\label{tab:steps}
\begin{center}
\small
\begin{tabular}{llccccccc}
\toprule
 & & \multicolumn{4}{c}{value still correct at position} & all steps & \multicolumn{2}{c}{final answer} \\
\cmidrule(lr){3-6}\cmidrule(lr){8-9}
Notation & Share & 2 & 3 & 4 & 5 & correct & steps ok & steps wrong \\
\midrule
none    & 0\%  & $100.0$ & $98.1$ & $92.1$ & $85.5$ & $84.9$ & $100.0$ & $3.5$ \\
formal  & 1\%  & $100.0$ & $98.7$ & $95.2$ & $92.8$ & $92.4$ & $99.9$ & $5.3$ \\
formal  & 5\%  & $100.0$ & $98.7$ & $96.5$ & $95.9$ & $95.5$ & $100.0$ & $8.8$ \\
formal  & 10\% & $100.0$ & $99.3$ & $97.7$ & $\mathbf{97.6}$ & $\mathbf{97.5}$ & $100.0$ & $5.3$ \\
English & 1\%  & $100.0$ & $97.9$ & $91.9$ & $84.1$ & $83.6$ & $100.0$ & $3.3$ \\
English & 5\%  & $100.0$ & $97.3$ & $91.9$ & $84.5$ & $84.1$ & $100.0$ & $2.5$ \\
English & 10\% & $100.0$ & $98.4$ & $93.2$ & $87.5$ & $87.1$ & $100.0$ & $3.1$ \\
\bottomrule
\end{tabular}
\end{center}
\end{table}


\subsection{FOLIO by gold label}
\label{sec:folio-table}

\begin{table}[H]
\caption{FOLIO by gold label, accuracy in percent. 35.5 percent of the evaluation items are true and the baseline answers ``true'' on 51.7 percent.}
\label{tab:folio}
\begin{center}
\begin{tabular}{llcccc}
\toprule
Notation & Share & false & true & uncertain & overall \\
\midrule
none    & 0\%  & $32.3$ & $80.1$ & $38.6$ & $51.4$ \\
formal  & 1\%  & $32.3$ & $80.6$ & $43.0$ & $53.0$ \\
formal  & 5\%  & $\mathbf{45.2}$ & $80.1$ & $45.9$ & $\mathbf{57.8}$ \\
formal  & 10\% & $37.1$ & $79.2$ & $\mathbf{52.2}$ & $57.1$ \\
formal  & 25\% & $40.3$ & $76.4$ & $49.3$ & $56.2$ \\
English & 1\%  & $32.3$ & $83.8$ & $42.5$ & $54.0$ \\
English & 5\%  & $44.1$ & $78.7$ & $43.5$ & $56.2$ \\
English & 10\% & $42.5$ & $75.5$ & $44.9$ & $55.0$ \\
\bottomrule
\end{tabular}
\end{center}
\end{table}


\subsection{Multi-hop question answering}
\label{sec:multihop-table}

\begin{table}[t]
\caption{Multi-hop question answering, F1 in percent, mean over seeds. The
three benchmarks retrieve from supplied passages and ask a question that
requires two or three hops. The bottom row gives the range over the three
control seeds. Gains are present and small, and the English rendering at ten
percent is the only condition above that range on all three.}
\label{tab:multihop}
\begin{center}
\begin{tabular}{llccc}
\toprule
Notation & Share & HotpotQA & 2WikiMultihopQA & MuSiQue \\
\midrule
none    & 0\%  & $55.5$ & $39.7$ & $31.1$ \\
formal  & 1\%  & $54.6$ & $39.7$ & $31.3$ \\
formal  & 5\%  & $56.1$ & $41.3$ & $32.7$ \\
formal  & 10\% & $54.8$ & $41.3$ & $31.0$ \\
formal  & 25\% & $55.2$ & $40.3$ & $31.6$ \\
English & 1\%  & $55.4$ & $40.0$ & $32.2$ \\
English & 5\%  & $55.4$ & $41.2$ & $\mathbf{34.2}$ \\
English & 10\% & $\mathbf{56.6}$ & $\mathbf{41.6}$ & $32.8$ \\
\midrule
\multicolumn{2}{l}{control seed range} & $0.1$ & $1.1$ & $1.5$ \\
\bottomrule
\end{tabular}
\end{center}
\end{table}


\subsection{Derivations without instruction data}
\label{sec:alone-table}

Table~\ref{tab:alone} gives the ablation of Section~\ref{sec:ablations} that
removes the instruction data.

\begin{table}[H]
\caption{Training on derivations without instruction data. Asked a benchmark question, these models emit the opening token of the derivation format and stop.}
\label{tab:alone}
\begin{center}
\small
\begin{tabular}{lcccc}
\toprule
training data & \branchproof{} & PW d3 & PW d5 & FOLIO \\
\midrule
instruction data only & $0.3$ & $39.3$ & $41.2$ & $51.4$ \\
5 percent \formal{} in instruction data & $99.3$ & $50.0$ & $48.9$ & $57.8$ \\
\formal{} derivations only & $100.0$ & $0.0$ & $0.0$ & $0.0$ \\
\nat{} derivations only & $100.0$ & $0.0$ & $0.0$ & $0.0$ \\
\bottomrule
\end{tabular}
\end{center}
\end{table}


\subsection{Both renderings in one mixture}
\label{sec:mixed-table}

Every mixture in the main text carries a single rendering of the derivations.
Table~\ref{tab:mixed} reports mixtures that carry both renderings of the same
derivations, in two forms. The plain form interleaves the two renderings with
no marker of which one a given derivation uses. The conditioned form prepends
one sentence to each training prompt that names the notation the derivation
will use, and we evaluate it with the same two sentences and with no sentence
at all. We train the plain form at shares of 2, 5, 10 and 20 percent and the
conditioned form at 2, 5, 10 and 20 percent, on Dolci and on the same base model as
Table~\ref{tab:corpus}.

Transfer follows the share and not the conditioning. The plain mixture at 10
percent gains 13.3 points on ProofWriter at inference depth 3, 9.8 points at
depth 5 and 5.3 points on FOLIO over the control, which is at the level of the
5 percent \formal{} mixture, while the plain mixtures at 2 and 5 percent gain
less and are not ordered by their share. The plain mixture at 20 percent gains
10.3 points at inference depth 3 and 1.8 points on FOLIO, below the 10 percent
mixture on both. At the matched share of 5 percent the conditioned mixture is
2.4 points above the plain one at inference depth 3 and 1.5 points below it on
FOLIO, at the matched share of 10 percent it is 4.0 points below it at
inference depth 3 and 2.0 points below it on FOLIO, and at the matched share of
20 percent it is 2.4 points below it at inference depth 3 and 0.5 points above
it on FOLIO.

What the conditioning buys is control of the notation. The conditioned models
write a derivation the checker accepts under either instruction, at 96.5 to
100 percent in both renderings and at all four shares. At 2 and 5 percent
they write almost no accepted derivation when neither instruction is present,
at 20.0 and 15.0 percent, while at 10 and 20 percent they write one at 99.0 and
100.0 percent, so the instruction is required for the notation at the larger
shares and not for the derivation itself. The plain model at 2 percent writes an accepted derivation when it is
left to choose, and it also follows the formal instruction in form, with a
formal block in 100 percent of its answers at chain length 25, but the checker
accepts 1.5 percent of those derivations. Training on both renderings without a
marker therefore leaves a model that can imitate the notation it is asked for
without being able to derive in it, which the one added sentence repairs. The
failure is confined to the smallest share, and the plain mixtures at 5, 10 and
20 percent derive under the formal instruction at 96.5 percent and above.

\begin{table}[t]
\caption{Mixtures that contain both renderings of the same derivations, with
and without a notation instruction in training. The first three columns give
the percentage of held-out generator problems at chain length 25 for which the
checker accepts the written derivation, under a prompt that asks for no
particular notation, a prompt that asks for a formal derivation and a prompt
that asks for an English one. PW d3 and PW d5 are accuracy in percent on
ProofWriter at inference depths 3 and 5 and FOLIO is accuracy in percent. The
first two rows repeat the control and the 5 percent \formal{} mixture of
Table~\ref{tab:corpus} for comparison. All rows are single runs. A dash marks a
measurement we did not run.}
\label{tab:mixed}
\begin{center}
\footnotesize
\begin{tabular}{lcccccc}
\toprule
 & \multicolumn{3}{c}{\branchproof{} by prompt} & & & \\
\cmidrule(lr){2-4}
Mixture & none & formal & English & PW d3 & PW d5 & FOLIO \\
\midrule
control                        & $0.3$  & --     & --      & $39.3$ & $41.2$ & $51.4$ \\
5\% \formal{}                  & $99.3$ & --     & --      & $50.0$ & $48.9$ & $57.8$ \\
\midrule
2\% both renderings            & $92.5$ & $1.5$  & $99.0$  & $47.0$ & $47.0$ & $51.2$ \\
5\% both renderings            & $95.0$ & $96.5$ & $100.0$ & $40.8$ & $42.2$ & $54.2$ \\
10\% both renderings           & $99.5$ & $98.5$ & $99.5$  & $52.6$ & $51.0$ & $56.7$ \\
20\% both renderings           & $100.0$ & --    & --      & $49.6$ & $49.4$ & $53.2$ \\
2\% both, conditioned          & $20.0$ & $97.0$ & $99.5$  & $47.0$ & $46.8$ & $50.7$ \\
5\% both, conditioned          & $15.0$ & $96.5$ & $100.0$ & $43.2$ & $43.6$ & $52.7$ \\
\bottomrule
\end{tabular}
\end{center}
\end{table}


\subsection{ProofWriter by inference depth}
\label{sec:depth-figure}

% ProofWriter accuracy by inference depth, means over seeds. Numbers match Table 1.
\begin{figure}[H]
\centering
\begin{tikzpicture}
\begin{axis}[
  width=0.62\linewidth, height=5.2cm,
  xmin=-0.3, xmax=4.3, ymin=28, ymax=68,
  xtick={0,1,2,3,4}, xticklabels={0,1,2,3,5},
  ytick={30,40,50,60},
  tick align=outside, tick pos=left,
  xlabel={inference depth of the ProofWriter item},
  ylabel={accuracy (\%)},
  grid=major, grid style={gray!16, line width=0.3pt},
  axis line style={gray!55},
  label style={font=\footnotesize}, tick label style={font=\scriptsize},
  legend style={at={(1.03,0.5)}, anchor=west, draw=none, fill=none,
                font=\scriptsize, row sep=-1pt},
  every axis plot/.append style={line width=1.0pt},
]
  \addplot [gray, mark=o, mark size=1.5pt, dashed] coordinates {(0,50.6) (1,30.9) (2,39.1) (3,39.3) (4,41.2)};
  \addplot [formalcol, mark=*, mark size=1.5pt] coordinates {(0,53.6) (1,36.2) (2,49.1) (3,50.0) (4,48.9)};
  \addplot [natcol, mark=square*, mark size=1.5pt] coordinates {(0,54.3) (1,37.3) (2,50.4) (3,53.2) (4,51.4)};
  \addplot [formalcol, mark=triangle*, mark size=2pt, densely dotted] coordinates {(0,63.6) (1,42.6) (2,51.0) (3,51.6) (4,49.8)};
  \legend{control, Formal 5\%, English 10\%, Formal 25\%}
\end{axis}
\end{tikzpicture}
\caption{Accuracy on ProofWriter for Qwen2.5-7B by the inference depth of
the item, in percent, mean over runs.}
\label{fig:depth}
\end{figure}


\subsection{Log-likelihood probe}
\label{sec:probe-table}

\begin{table}[H]
\caption{Ranking of decidable ProofWriter items by the log-likelihood margin
between ``true'' and ``false'' on Qwen2.5-7B, area under the ranking curve
(AUC) and the probability of preferring ``true'', both in percent, for 1 run with a second run in parentheses. The pooled column mixes the
two polarities and the next two columns rank within a polarity. \longdoc{} is a length-matched control that
replaces the same share of the mixture with ordinary long documents.}
\label{tab:probe}
\begin{center}
\small
\begin{tabular}{lcccc}
\toprule
Mixture & AUC pooled & AUC negated & AUC positive & $P(\text{argmax}=\text{true})$ \\
\midrule
\ctrl{}    & $66.2$ ($65.4$) & $87.2$ ($86.9$) & $86.0$ ($85.8$) & $64.0$ ($63.0$) \\
\longdoc{} & $66.3$ ($66.4$) & $87.6$ ($87.7$) & $86.7$ ($86.7$) & $64.0$ ($64.0$) \\
\formal{}  & $71.1$ ($70.7$) & $88.9$ ($88.9$) & $86.8$ ($86.8$) & $55.0$ ($54.0$) \\
\nat{}     & $77.0$          & $89.8$          & $87.3$          & $48.0$ \\
\bottomrule
\end{tabular}
\end{center}
\end{table}

\subsection{Perturbation probes}
\label{sec:perturbation}

The first probe toggles the polarity of the claim and leaves the theory
untouched, so a model that ignores negation answers both versions
identically. The control does so on 37.6 percent of decidable items and the
length-matched control on 37.6 percent, against 28.5 percent for the formal
and 24.2 percent for the English mixture, with a second run within half a
point throughout. Every model answers the positive version of a negated and
false item correctly about 68 percent of the time, so the models separate on
the negated version alone.

The second probe removes the premises that support the claim, which makes
the gold label unknown, and asks whether the additional ``false'' answers
survive. They largely do not, which is what a premise-grounded answer
predicts. The probability that the formal mixture answers ``false'' on these
items falls from 53.0 to 27.0 percent once the supporting facts are gone, and the
control falls from 35.0 to 20.0 percent. The English mixture retains more, falling
from 67.0 to 46.0 percent, so about two thirds of its advantage on this cell
survives the removal of the evidence for it. The formal rendering therefore
buys a smaller gain that is more firmly tied to the premises, and the
English rendering a larger gain of which part is a shifted default for
negated claims.

\section{Qualitative Examples}
\label{sec:examples}

Each example below is an item that the Qwen2.5-7B control answers
incorrectly and that the models trained with 10 percent formal and 10
percent English derivations both answer correctly, all from the same run
of each condition. The
responses are reproduced verbatim apart from the omission of repeated
prefixes.

\paragraph{ProofWriter.} The theory is ``The bear needs the dog. The cow is
kind. The cow is nice. The cow likes the dog. The cow needs the dog. The dog
chases the rabbit. The dog likes the bear. The dog likes the rabbit. The
rabbit chases the bear. The rabbit chases the cow. The rabbit is blue. The
rabbit is kind. The rabbit likes the cow. The rabbit needs the dog. If
someone needs the rabbit then the rabbit likes the bear. If someone needs the
rabbit and they chase the dog then they need the bear. If someone is kind
then they like the dog. If someone likes the bear and they are blue then
they chase the rabbit. If the rabbit likes the dog then the dog is kind. If
someone chases the dog then the dog likes the cow. If someone likes the dog
then they are round. If the rabbit likes the dog and the rabbit chases the
cow then the dog likes the cow. If someone needs the dog and the dog is kind
then the dog needs the rabbit.'' The claim is ``The dog does not need the
rabbit'', which is false, since the rabbit is kind, therefore likes the dog,
therefore the dog is kind, and the rabbit needs the dog, so the dog needs
the rabbit. The control answers ``True'' and both trained models answer
``False''. The item in Section~\ref{sec:prior} has the same structure.

\paragraph{FOLIO.} The premises are ``A Japanese game company created the
game the Legend of Zelda. All games on the Top 10 list are made by Japanese
game companies. If a game sells more than one million copies, then it will
be included in the Top 10 list. The Legend of Zelda sold more than one
million copies.'' The claim ``The Legend of Zelda is not on the Top 10 list''
is false. The control answers ``Uncertain'' and both trained models answer
``False''. A second item has the premises ``Beethoven leads the Vienna Music
Society. Vienna Music Society is an orchestra. Orchestras are led by
conductors.'' and the claim ``Beethoven is not a conductor''. The control
answers ``True'' and both trained models answer ``False''.

\paragraph{Web of lies.} A second item beyond the one in
Section~\ref{sec:procedure}.

\begin{small}
\begin{quote}
\textbf{Question.} Jim lies. Yoland says Jim lies. Willian says Yoland tells
the truth. Delbert says Willian lies. Shalonda says Delbert tells the truth.
Does Shalonda tell the truth?

\textbf{Control.} (1) Jim lies. (2) Since Jim lies and Yoland says Jim lies,
Yoland tells the truth. (3) Willian tells the truth. (4) Since Willian tells
the truth and Delbert says Willian lies, Delbert lies. (5) Since we know from
(4) that Delbert \emph{tells the truth}, and Shalonda says Delbert tells the
truth, Shalonda tells the truth. So the answer is Yes.

\textbf{With derivations.} Steps (1) to (4) as above. (5) Since we know from
(4) that Delbert lies, and Shalonda says Delbert tells the truth, Shalonda
lies. So the answer is No.
\end{quote}
\end{small}

\paragraph{Tracking shuffled objects.} Alice, Bob, Claire, Dave and Eve
hold a brown, purple, orange, pink and blue present respectively, and the
pairs Claire and Eve, Claire and Bob, Alice and Dave, Bob and Claire, and
Bob and Alice swap in turn. The question is which present Eve holds at the
end. Both models write out the full state after every swap. The control
performs the second swap between Claire and Bob incorrectly, giving Bob the
orange present and returning the blue present to Eve, and answers blue,
whereas the trained models carry the state correctly through all five swaps
and answer orange. On a second item the control reaches the correct final
state and then selects the wrong option letter, which the trained models do
not.

\paragraph{GPQA-Diamond.} An uncharged spherical conductor of radius $R$
contains an off-centre cavity with a charge $+q$ inside, and the question
asks for the field at a point $P$ outside the conductor at distance $L$ from
the centre of the conductor and $l$ from the centre of the cavity. The
control writes that the field ``is the same as if the charge were
concentrated at the center of the cavity'' and answers $q/(4\pi\epsilon_0
l^2)$. The trained models write that the conductor shields the position of
the charge, so the field is that of a charge at the centre of the conductor,
and answer $q/(4\pi\epsilon_0 L^2)$. In a second item the question gives a
size estimate for a loop diagram containing the factor $1/(4\pi)^6$ and asks
for the number of loops. The control reads the exponent as the loop count and
answers 6, whereas the model trained with formal derivations writes that each
loop contributes a factor $1/(4\pi)^2$ and answers 3.

\subsection{Web of lies}
\label{sec:wol-example}

The item below shows the failure of the Qwen2.5-7B control on web of lies
described in Section~\ref{sec:procedure} and its absence after training.

\begin{small}
\begin{quote}
\textbf{Question.} Vina tells the truth. Helene says Vina lies. Kandi says
Helene tells the truth. Jamey says Kandi lies. Ka says Jamey lies. Does Ka
tell the truth?

\textbf{Control.} (1) Vina tells the truth. (2) Since Vina tells the truth
and Helene says Vina lies, Helene lies. (3) Since we know from (2) that
Helene \emph{tells the truth}, and Kandi says Helene tells the truth, Kandi
tells the truth. (4) Jamey lies. (5) Ka tells the truth. So the answer is
Yes.

\textbf{With 10 percent formal derivations.} (1) Vina tells the truth. (2)
Helene lies. (3) Since we know from (2) that Helene lies, and Kandi says
Helene tells the truth, Kandi lies. (4) Jamey tells the truth. (5) Ka lies.
So the answer is No.
\end{quote}
\end{small}

\section{Verification of Derivations}
\label{sec:verification}

A derivation in the trained format is
checked by the same program that wrote the training data. On the generator's
own problems, a derivation the checker accepts carries the correct answer in
99 percent of samples, so selecting the first accepted derivation out of
16 matches the accuracy of selecting with the gold answer. On
ProofWriter the checker reproduces the published label on 84 to 99 percent
of items depending on the depth, and an accepted derivation is correct
between 53 and 69 percent of the time, because a third of the items ask
whether a claim is underivable and a forward chain cannot show that. The
coupling between validity and correctness therefore holds where the checker
is complete for the task, which makes the trained format a candidate for
verifier-based training on deduction \citep{lambert2025tulu, yue2025does},
although we do not pursue this here.

\section{Notation Crossover on the Generator}
\label{sec:crossover}

A controlled study on the generator alone shows that the ordering of the two
renderings in depth generalisation depends on the depth of the training
proofs. We fine-tune OLMo-3-7B with LoRA adapters on 50{,}000 paired
examples per condition, one condition per rendering, for 10{,}000 optimizer
steps, and we evaluate at chain lengths beyond the longest one seen in
training with 16 samples per prompt and 1 greedy decode.
Figure~\ref{fig:crossover} gives the resulting curves. English supervision
generalises further as long as the training proofs reach at most 20 hops. At
25 hops the ordering reverses, and formal supervision reaches $75.0\pm9.5$
percent out-of-distribution accuracy against $17.8\pm4.7$ percent for English
supervision. Sixteen samples move the reversal to 15 hops, so formal
supervision first widens the set of problems the model can reach and only
later improves its first attempt.

% Generated from the three-seed grid. Shaded regions are one standard
% deviation over seeds. Do not edit by hand.
\pgfplotsset{crossoverstyle/.style={
    width=0.34\linewidth, height=4.5cm,
    xmin=4, xmax=26, ymin=0, ymax=105,
    xtick={5,10,15,20,25}, ytick={0,25,50,75,100},
    tick align=outside, tick pos=left,
    xlabel={maximum training depth},
    grid=major, grid style={gray!16, line width=0.3pt},
    axis line style={gray!55},
    label style={font=\footnotesize}, tick label style={font=\scriptsize},
    title style={font=\small, yshift=-2pt},
    every axis plot/.append style={line width=1.0pt},
}}
\begin{figure}[H]
\centering
\begin{tikzpicture}
\begin{axis}[crossoverstyle, name=plot1,  title={greedy}, ylabel={out-of-distribution accuracy (\%)},]
  \addplot [name path=logicu, draw=none, forget plot] coordinates {(5,16.7) (10,35.3) (15,31.6) (20,33.4) (25,89.0)};
  \addplot [name path=logicl, draw=none, forget plot] coordinates {(5,5.0) (10,30.9) (15,24.6) (20,10.7) (25,75.1)};
  \addplot [formalcol, opacity=0.16, forget plot] fill between [of=logicu and logicl];
  \addplot [name path=nlu, draw=none, forget plot] coordinates {(5,64.3) (10,68.4) (15,62.4) (20,88.2) (25,70.5)};
  \addplot [name path=nll, draw=none, forget plot] coordinates {(5,48.4) (10,51.6) (15,52.5) (20,48.6) (25,27.0)};
  \addplot [natcol, opacity=0.16, forget plot] fill between [of=nlu and nll];
  \addplot [formalcol, mark=*, mark size=1.4pt] coordinates {(5,10.8) (10,33.1) (15,28.1) (20,22.1) (25,82.1)};
  \addplot [natcol, mark=*, mark size=1.4pt] coordinates {(5,56.3) (10,60.0) (15,57.4) (20,68.4) (25,48.8)};
\end{axis}
\begin{axis}[crossoverstyle, name=plot2, legend style={at={(0.5,-0.32)}, anchor=north, legend columns=2, draw=none, fill=none, font=\footnotesize, /tikz/every even column/.append style={column sep=8pt}}, at={($(plot1.east)+(0.9cm,0)$)}, anchor=west, title={sampled \passk{1}}, yticklabels={},]
  \addplot [name path=logicu, draw=none, forget plot] coordinates {(5,14.6) (10,32.6) (15,29.0) (20,34.6) (25,84.5)};
  \addplot [name path=logicl, draw=none, forget plot] coordinates {(5,5.0) (10,29.1) (15,23.1) (20,11.9) (25,65.5)};
  \addplot [formalcol, opacity=0.16, forget plot] fill between [of=logicu and logicl];
  \addplot [name path=nlu, draw=none, forget plot] coordinates {(5,59.1) (10,66.1) (15,50.5) (20,36.6) (25,22.5)};
  \addplot [name path=nll, draw=none, forget plot] coordinates {(5,45.7) (10,49.8) (15,48.5) (20,25.8) (25,13.1)};
  \addplot [natcol, opacity=0.16, forget plot] fill between [of=nlu and nll];
  \addplot [formalcol, mark=*, mark size=1.4pt] coordinates {(5,9.8) (10,30.8) (15,26.1) (20,23.2) (25,75.0)};
  \addplot [natcol, mark=*, mark size=1.4pt] coordinates {(5,52.4) (10,58.0) (15,49.5) (20,31.2) (25,17.8)};
  \legend{Formal, English}
\end{axis}
\begin{axis}[crossoverstyle, name=plot3, at={($(plot2.east)+(0.9cm,0)$)}, anchor=west, title={sampled \passk{16}}, yticklabels={},]
  \addplot [name path=logicu, draw=none, forget plot] coordinates {(5,65.7) (10,70.8) (15,76.5) (20,92.9) (25,99.9)};
  \addplot [name path=logicl, draw=none, forget plot] coordinates {(5,34.3) (10,64.9) (15,63.3) (20,67.9) (25,99.3)};
  \addplot [formalcol, opacity=0.16, forget plot] fill between [of=logicu and logicl];
  \addplot [name path=nlu, draw=none, forget plot] coordinates {(5,83.1) (10,81.9) (15,62.8) (20,61.9) (25,48.3)};
  \addplot [name path=nll, draw=none, forget plot] coordinates {(5,63.9) (10,60.4) (15,51.0) (20,28.7) (25,36.7)};
  \addplot [natcol, opacity=0.16, forget plot] fill between [of=nlu and nll];
  \addplot [formalcol, mark=*, mark size=1.4pt] coordinates {(5,50.0) (10,67.8) (15,69.9) (20,80.4) (25,99.6)};
  \addplot [natcol, mark=*, mark size=1.4pt] coordinates {(5,73.5) (10,71.2) (15,56.9) (20,45.3) (25,42.5)};
\end{axis}
\end{tikzpicture}
\caption{Out-of-distribution answer accuracy on \branchproof{} against maximum
training depth, for three decoding rules. Shaded regions give one standard
deviation over 3 runs. Greedy decoding and \passk{1} place the crossover
at maximum training depth 25. Sampling with sixteen draws places it at 15, ten
depth levels earlier.}
\label{fig:crossover}
\end{figure}


The crossover follows from the supervision alone, since the prompt, the
answer and the latent proof are identical in the two conditions at every
depth. It survives a control that equates the number of supervised target
tokens by subsampling the longer condition, a control that removes a schema
marker which identifies the active branch on 80 percent of training examples
(out-of-distribution \passk{1} of $43.1\pm7.7$ percent for \formal{} and
$78.5\pm14.9$ for \nat{}), and a replication with a second run. Concatenating
both renderings into one target fails, with $0.5\pm0.4$ percent
out-of-distribution \passk{1}, whereas placing an output-mode tag on separate
examples reaches $50.3\pm6.4$ percent in formal mode at 7B and $85.0\pm8.7$
at 32B.

\section{Midtraining}
\label{sec:midtrain-detail}

The same derivations can enter during midtraining. We midtrained Qwen2.5-7B
on 2.5 billion tokens, 2{,}385 steps at sequence length 8{,}192 with a
global batch of 128 sequences, with a tenth of the loss tokens replaced by
72{,}000 generated documents packed without splitting, and then applied the
same instruction tuning to every condition. The conditions are the control,
a length-matched control that replaces the same tokens by ordinary long
documents (\longdoc{}), and the formal and English renderings, with 1
midtraining run per condition. Table~\ref{tab:midtrain} reports the
results. That route also produces the gain on ProofWriter, 12.6 points for
the English rendering at inference depth 3, which a second instruction-tuning
run reproduces to within 0.1 points, and the gain again lies in the items
whose gold label is false. It also leaves the general benchmarks flat apart
from a loss of 2 to 3 points on LogiQA that the second run reproduces.
Midtraining is the more expensive route and buys no more, since the
instruction-tuning route reaches the same accuracy, generalises further on
long chains, carries no cost on LogiQA, and adds nothing when the two are
applied in sequence. The midtrained base models write checker-valid
derivations on every generator problem up to chain length 25 when prompted
in the format of their training documents, and the control base writes
none. Ordinary instruction tuning removes this ability, and the lower half
of Table~\ref{tab:midtrain} shows what remains of it after instruction
tuning, where the derivations raise greedy accuracy on the generator's
problems by 20 to 50 points and \passk{16} by 20 to 40 points.

\begin{table}[H]
\caption{Midtraining on Qwen2.5-7B followed by instruction tuning without
derivations, in percent, 1 run per condition. The upper part gives
ProofWriter accuracy at each inference depth on 500 items per depth, and
the accuracy on the items pooled over depths 1 to 5 by gold label, together
with the fraction of pooled items on which the model answers ``true''. The
lower part gives accuracy on held-out generator problems at chain lengths 5
to 25, 200 items per length, under greedy decoding and as \passk{16} at
temperature 0.8.}
\label{tab:midtrain}
\begin{center}
\small
\setlength{\tabcolsep}{4pt}
\begin{tabular}{lccccccccc}
\toprule
 & \multicolumn{5}{c}{ProofWriter by inference depth} & \multicolumn{4}{c}{pooled depths 1 to 5} \\
\cmidrule(lr){2-6}\cmidrule(lr){7-10}
Condition & 0 & 1 & 2 & 3 & 5 & answers true & gold true & gold false & two-class \\
\midrule
\ctrl{}     & 44.2 & 32.8 & 44.0 & 44.4 & 46.0 & 66.3 & 70.4 & 38.5 & 54.4 \\
\longdoc{}  & 45.8 & 32.2 & 44.4 & 45.4 & 44.0 & 65.9 & 69.5 & 38.6 & 54.0 \\
\formal{}   & 48.0 & 35.8 & 49.4 & 50.0 & 50.2 & 56.4 & 66.5 & 54.2 & 60.3 \\
\nat{}      & \textbf{48.8} & \textbf{38.8} & \textbf{56.2} & \textbf{57.0} & \textbf{56.4} & 49.5 & 68.3 & \textbf{68.0} & \textbf{68.2} \\
\midrule
 & & \multicolumn{5}{c}{generator problems by chain length} \\
\cmidrule(lr){3-7}
Decoding & Condition & 5 & 10 & 15 & 20 & 25 \\
\midrule
greedy & \ctrl{} & 5.0 & 3.5 & 11.0 & 9.0 & 6.5 \\
greedy & \longdoc{} & 7.5 & 11.0 & 12.0 & 15.0 & 12.0 \\
greedy & \formal{} & 31.5 & 17.5 & 27.5 & 21.5 & 19.5 \\
greedy & \nat{} & \textbf{52.5} & \textbf{34.5} & \textbf{30.5} & \textbf{24.5} & \textbf{29.5} \\
\passk{16} & \ctrl{} & 58.0 & 55.5 & 58.5 & 60.0 & 53.0 \\
\passk{16} & \longdoc{} & 63.0 & 63.5 & 61.0 & 63.5 & 56.0 \\
\passk{16} & \formal{} & \textbf{97.0} & 85.5 & 79.0 & 74.0 & 71.5 \\
\passk{16} & \nat{} & 95.5 & \textbf{91.5} & \textbf{84.5} & \textbf{78.0} & \textbf{82.5} \\
\bottomrule
\end{tabular}
\end{center}
\end{table}
